\section{IMU Sensor Calibration \& System Identification}
    센서 캘리브레이션(Sensor Calibration)은 IMU의 측정값에 포함된 체계적 오차(systematic error)를 정량화하고 보정하는 과정이다. 2장에서 정의한 측정 모델에는 바이어스(bias), 스케일 팩터 오차(scale factor error), 축 정렬 오차(misalignment error) 등의 결정론적 오차가 포함되어 있으며, 이를 보정하지 않으면 자세 추정 성능이 크게 저하된다. 본 장에서는 가속도계, 자이로스코프, 지자기계 각각의 캘리브레이션 방법을 기술하고, 캘리브레이션 결과의 검증 방법을 제시한다.\\

    \subsection{Accelerometer Calibration}
        가속도계의 실제 측정 모델은 바이어스, 스케일 팩터, 축 정렬 오차를 포함하여 다음과 같이 일반화할 수 있다:\\

        \begin{equation}
            \tilde{\mathbf{a}} = \mathbf{T}_a \mathbf{S}_a \left( \mathbf{a}_{\text{true}} - \mathbf{b}_a \right) + \mathbf{n}_a
            \label{eq:accel_full_model}
        \end{equation}

        \noindent 여기서 $\mathbf{T}_a \in \mathbb{R}^{3\times 3}$은 축 정렬 오차 행렬(misalignment matrix), $\mathbf{S}_a = \text{diag}(s_{ax},\, s_{ay},\, s_{az})$는 스케일 팩터 행렬, $\mathbf{b}_a \in \mathbb{R}^3$은 바이어스 벡터이다.\\

        이를 하나의 보정 행렬로 통합하면:\\

        \begin{equation}
            \mathbf{a}_{\text{true}} = \mathbf{K}_a^{-1} \tilde{\mathbf{a}} + \mathbf{b}_a, \qquad \mathbf{K}_a = \mathbf{T}_a \mathbf{S}_a
            \label{eq:accel_calib_matrix}
        \end{equation}

        \subsubsection{Bias and Scale Factor Estimation}
            \paragraph{6-위치 정적 캘리브레이션}

                가속도계 캘리브레이션의 가장 일반적인 방법은 6-위치 정적 캘리브레이션(6-position static calibration)이다. 센서를 각 축의 양/음 방향이 중력 방향과 일치하도록 6개의 위치에 순서대로 정치시키고, 각 위치에서의 평균 측정값을 수집한다.\\

                이상적인 경우, 각 위치에서 측정된 비력의 크기는 중력 가속도 $g$와 같아야 한다:\\

                \begin{equation}
                    \left\| \mathbf{K}_a^{-1}(\tilde{\mathbf{a}}_i - \mathbf{b}_a) \right\| = g, \qquad i = 1, \ldots, 6
                    \label{eq:accel_6pos_constraint}
                \end{equation}

                6개의 측정 벡터를 행렬로 표현하면:\\

                \begin{equation}
                    \mathbf{A} = \begin{bmatrix}
                        \tilde{a}_{x,1} & \tilde{a}_{y,1} & \tilde{a}_{z,1} \\
                        \vdots & \vdots & \vdots \\
                        \tilde{a}_{x,6} & \tilde{a}_{y,6} & \tilde{a}_{z,6}
                    \end{bmatrix}
                    \label{eq:accel_meas_matrix}
                \end{equation}

                각 축의 바이어스와 스케일 팩터는 양/음 방향 쌍으로부터 다음과 같이 추정한다:\\

                \begin{align}
                    b_{a,i}   &= \frac{\tilde{a}_{i,+} + \tilde{a}_{i,-}}{2} \label{eq:accel_bias_est} \\
                    s_{a,i}   &= \frac{\tilde{a}_{i,+} - \tilde{a}_{i,-}}{2g} \label{eq:accel_scale_est}
                \end{align}

                \noindent 여기서 $\tilde{a}_{i,+}$, $\tilde{a}_{i,-}$는 각각 $i$축 양/음 방향 정치 시의 평균 측정값이다 ($i \in \{x, y, z\}$).\\

            \paragraph{최소제곱 추정}
                보다 정확한 추정을 위해 최소제곱법(Least Squares)을 적용할 수 있다. 파라미터 벡터 $\boldsymbol{\theta} = [\mathbf{b}_a^\top,\, \mathbf{s}_a^\top]^\top \in \mathbb{R}^6$에 대한 선형 관측 방정식은 다음과 같다:\\

                \begin{equation}
                    \mathbf{y} = \boldsymbol{\Phi}\boldsymbol{\theta} + \boldsymbol{\epsilon}
                    \label{eq:accel_ls}
                \end{equation}

                최소제곱 추정량은 다음과 같이 주어진다:\\

                \begin{equation}
                    \hat{\boldsymbol{\theta}} = \left(\boldsymbol{\Phi}^\top \boldsymbol{\Phi}\right)^{-1} \boldsymbol{\Phi}^\top \mathbf{y}
                    \label{eq:accel_ls_solution}
                \end{equation}

        \subsubsection{Misalignment Correction}
            축 정렬 오차는 센서의 감지 축이 이상적인 직교 좌표계 축과 미세하게 어긋나는 현상으로, 제조 공정의 한계에 기인한다. 축 정렬 오차 행렬 $\mathbf{T}_a$는 다음과 같이 표현된다:\\

            \begin{equation}
                \mathbf{T}_a = \begin{bmatrix}
                    1       & -\alpha_{xy} & \alpha_{xz} \\
                    \alpha_{yx} & 1       & -\alpha_{yz} \\
                    -\alpha_{zx} & \alpha_{zy} & 1
                \end{bmatrix}
                \label{eq:misalignment_matrix}
            \end{equation}

            \noindent 여기서 $\alpha_{ij}$는 $j$축에 대한 $i$축의 정렬 오차 각도(radian)이다. 일반적으로 MEMS IMU에서 정렬 오차는 $0.1^\circ$ 미만이나, 고정밀 응용에서는 반드시 보정이 필요하다.\\

            통합 보정 절차를 정리하면 다음과 같다:\\

            \begin{equation}
                \mathbf{a}_{\text{calib}} = \mathbf{T}_a^{-1} \mathbf{S}_a^{-1} \left( \tilde{\mathbf{a}} - \mathbf{b}_a \right)
                \label{eq:accel_full_calib}
            \end{equation}

    \subsection{Gyroscope Calibration}
        자이로스코프 캘리브레이션의 목적은 정적 바이어스를 추정·제거하고, Allan Variance 분석을 통해 잡음 파라미터를 정량화하는 것이다. 동적 캘리브레이션(정밀 회전 테이블 사용)은 본 연구의 범위를 벗어나므로, 정적 조건에서의 캘리브레이션에 집중한다.\\

        \subsubsection{Static Bias Estimation}
            정지 상태에서 자이로스코프의 출력은 실제 각속도($\boldsymbol{\omega} = \mathbf{0}$)와 바이어스, 잡음의 합이므로:\\

            \begin{equation}
                \tilde{\boldsymbol{\omega}}(t) = \mathbf{b}_g + \mathbf{n}_g(t)
                \label{eq:gyro_static}
            \end{equation}

            충분히 긴 시간 $T$ 동안의 평균을 취하면 잡음 성분이 상쇄되어 바이어스를 추정할 수 있다:\\

            \begin{equation}
                \hat{\mathbf{b}}_g = \frac{1}{N} \sum_{k=1}^{N} \tilde{\boldsymbol{\omega}}_k
                \label{eq:gyro_bias_est}
            \end{equation}

            추정 바이어스의 표준 오차는 다음과 같다:\\

            \begin{equation}
                \text{SE}(\hat{\mathbf{b}}_g) = \frac{\sigma_g}{\sqrt{N}} = \frac{\sigma_g}{\sqrt{f_s T}}
                \label{eq:gyro_bias_se}
            \end{equation}

            \noindent 여기서 $f_s$는 샘플링 주파수, $T$는 수집 시간이다. 예를 들어 $f_s = 100\,\text{Hz}$, $T = 60\,\text{s}$, $\sigma_g = 0.01\,\text{deg/s}$이면 $\text{SE} \approx 0.00013\,\text{deg/s}$로 충분히 정밀한 추정이 가능하다.\\

            보정된 자이로스코프 출력은 다음과 같다:\\

            \begin{equation}
                \boldsymbol{\omega}_{\text{calib}} = \tilde{\boldsymbol{\omega}} - \hat{\mathbf{b}}_g
                \label{eq:gyro_calib}
            \end{equation}

            단, 정적 바이어스는 온도 변화 및 시간에 따라 드리프트하므로, 매 실험 시작 전에 새로 추정하는 것이 권장된다.\\

        % ------------------------------------------------------------
        \subsubsection{ARW and Bias Instability Estimation}
            ARW 및 Bias Instability 계수는 2.2.2절에서 정의한 Allan Variance 분석을 통해 추정한다. 실험 절차는 다음과 같다:\\

            \begin{enumerate}
                \item 센서를 진동이 없는 정지 환경에 최소 2시간 이상 정치
                \item $f_s$로 자이로스코프 출력 $\tilde{\omega}(t)$ 수집
                \item 적분 시간 $\tau = \Delta t, 2\Delta t, \ldots, T/2$에 대해
                    식 \eqref{eq:avar_estimate}으로 $\hat{\sigma}^2(\tau)$ 계산
                \item log-log 플롯에서 기울기 분석으로 각 잡음 계수 추출
            \end{enumerate}

            \paragraph{ARW 계수 추정}
                $\tau = 1\,\text{s}$에서의 Allan Deviation 값으로부터 ARW 계수 $N$을 읽는다:\\

                \begin{equation}
                    N = \sigma(\tau)\big|_{\tau=1\,\text{s}} \quad [\text{deg}/\sqrt{\text{hr}}]
                    \label{eq:arw_coeff}
                \end{equation}

                또는 기울기 $-1/2$ 직선을 피팅하여 $\tau = 1$에서의 절편으로 추정할 수 있다.\\

            \paragraph{EKF 파라미터 변환}
                Allan Variance 분석에서 얻은 계수를 EKF의 공정 잡음 공분산 $\mathbf{Q}$로 변환하는 과정은 다음과 같다.\\

                ARW 계수 $N\,[\text{deg}/\sqrt{\text{hr}}]$로부터 연속시간 각속도 잡음 PSD를 구하면:\\

                \begin{equation}
                    S_g = N^2 \quad [\text{deg}^2/\text{hr}]
                    \label{eq:arw_psd}
                \end{equation}

                이를 이산시간 분산으로 변환하면:\\

                \begin{equation}
                    \sigma_w^2 = \frac{S_g}{f_s} = \frac{N^2}{f_s} \quad [\text{rad}^2]
                    \label{eq:arw_discrete}
                \end{equation}

                마찬가지로 RRW 계수 $K\,[\text{deg}/\text{hr}^{3/2}]$로부터:

                \begin{equation}
                    \sigma_{rw}^2 = K^2 \cdot \Delta t \quad [\text{rad}^2/\text{s}]
                    \label{eq:rrw_discrete}
                \end{equation}

                이 두 값이 5장 EKF 설계에서 공정 잡음 공분산 행렬 $\mathbf{Q}$의 대각 성분으로 사용된다.\\

    \subsection{Magnetometer Calibration}
        지자기계 캘리브레이션의 목적은 2장에서 기술한 Hard Iron 바이어스 $\mathbf{b}_m$과 Soft Iron 왜곡 행렬 $\mathbf{W}$를 추정하여 보정하는 것이다. 이상적으로 보정된 지자기계 출력의 크기는 위치에 무관하게 일정해야 하므로, 3차원 공간에서 구(sphere)를 형성해야 한다. 실제 측정값은 Hard Iron에 의해 구의 중심이 이동하고, Soft Iron에 의해 구가 타원체(ellipsoid)로 변형된다.\\

        \subsubsection{Hard Iron Correction}
            Hard Iron 보정은 측정된 자기장 벡터의 오프셋을 제거하는 과정이다. 센서를 다양한 방향으로 회전시키면서 측정값을 수집하면, 이상적인 경우 반지름 $\|\mathbf{m}_{\text{true}}\|$인 구를 형성해야 하지만, Hard Iron 바이어스 $\mathbf{b}_m$에 의해 중심이 이동한 구가 관측된다.\\

            Hard Iron 바이어스는 수집된 측정값의 최대·최솟값으로부터 다음과 같이 추정할 수 있다:\\

            \begin{equation}
                \hat{b}_{m,i} = \frac{\max(\tilde{m}_i) + \min(\tilde{m}_i)}{2}, \qquad i \in \{x, y, z\}
                \label{eq:hard_iron_est}
            \end{equation}

            보정된 자기장 벡터:\\

            \begin{equation}
                \mathbf{m}_{\text{hi}} = \tilde{\mathbf{m}} - \hat{\mathbf{b}}_m
                \label{eq:hard_iron_calib}
            \end{equation}

        \subsubsection{Soft Iron Correction (Ellipsoid Fitting)}
            Soft Iron 보정은 Hard Iron 보정 후에도 남아있는 타원체 변형을 구로 복원하는 과정이다. 이를 위해 타원체 피팅(Ellipsoid Fitting) 기법을 사용한다.\\

            \paragraph{타원체 방정식}
                3차원 타원체는 다음 일반 이차형식(quadratic form)으로 표현된다:\\

                \begin{equation}
                    \mathbf{m}^\top \mathbf{A} \mathbf{m} + \mathbf{b}^\top \mathbf{m} + c = 0
                    \label{eq:ellipsoid_general}
                \end{equation}

                \noindent 여기서 $\mathbf{A} \in \mathbb{R}^{3\times 3}$은 대칭 양정치 행렬, $\mathbf{b} \in \mathbb{R}^3$, $c \in \mathbb{R}$은 타원체의 형상과 위치를 결정하는 파라미터이다.\\

            \paragraph{최소제곱 피팅}
                $N$개의 측정값 $\{\mathbf{m}_{\text{hi},k}\}_{k=1}^N$을 이용하여 파라미터 벡터 $\boldsymbol{\theta}$를 최소제곱법으로 추정한다.\\

                \begin{equation}
                    \boldsymbol{\theta} = [A_{xx},\, A_{yy},\, A_{zz},\, A_{xy},\, A_{xz},\, A_{yz},\, b_x,\, b_y,\, b_z,\, c]^\top
                    \label{eq:ellipsoid_params}
                \end{equation}

                설계 행렬 $\boldsymbol{\Phi} \in \mathbb{R}^{N \times 10}$을 다음과 같이 구성한다:\\

                \begin{equation}
                    \boldsymbol{\Phi}_k = \left[ m_x^2,\, m_y^2,\, m_z^2,\, 2m_xm_y,\, 2m_xm_z,\, 2m_ym_z,\, 2m_x,\, 2m_y,\, 2m_z,\, 1 \right]
                    \label{eq:ellipsoid_design}
                \end{equation}

                파라미터 추정:\\

                \begin{equation}
                    \hat{\boldsymbol{\theta}} = \left(\boldsymbol{\Phi}^\top \boldsymbol{\Phi}\right)^{-1} \boldsymbol{\Phi}^\top \mathbf{1}_N
                    \label{eq:ellipsoid_ls}
                \end{equation}

            \paragraph{보정 행렬 계산}
                추정된 타원체 파라미터로부터 Soft Iron 보정 행렬 $\mathbf{W}^{-1}$을 계산한다. 행렬 $\mathbf{A}$를 고유값 분해(eigendecomposition)하면:\\

                \begin{equation}
                    \mathbf{A} = \mathbf{V} \boldsymbol{\Lambda} \mathbf{V}^\top
                    \label{eq:eigen_decomp}
                \end{equation}

                \noindent 여기서 $\mathbf{V}$는 고유벡터 행렬, $\boldsymbol{\Lambda} = \text{diag}(\lambda_1, \lambda_2, \lambda_3)$는 고유값 대각 행렬이다.\\

                Soft Iron 보정 행렬은 다음과 같이 계산된다:\\

                \begin{equation}
                    \mathbf{W}^{-1} = \sqrt{\bar{\lambda}} \cdot \mathbf{V} \boldsymbol{\Lambda}^{-1/2} \mathbf{V}^\top, \qquad \bar{\lambda} = \frac{\lambda_1 + \lambda_2 + \lambda_3}{3}
                    \label{eq:soft_iron_matrix}
                \end{equation}

                최종 보정된 자기장 벡터:\\

                \begin{equation}
                    \mathbf{m}_{\text{calib}} = \mathbf{W}^{-1} \left( \tilde{\mathbf{m}} - \hat{\mathbf{b}}_m \right)
                    \label{eq:mag_full_calib}
                \end{equation}

                보정이 완료되면 $\mathbf{m}_{\text{calib}}$는 다양한 방향에서 일정한 크기를 가져야 한다:\\

                \begin{equation}
                    \left\| \mathbf{m}_{\text{calib}} \right\| \approx \left\| \mathbf{m}_{\text{true}} \right\| = \text{const.}
                    \label{eq:mag_calib_check}
                \end{equation}

    \subsection{Calibration Validation}
        캘리브레이션의 유효성은 보정 전후의 측정값을 정량적으로 비교함으로써 검증한다.\\
        
        \subsubsection{Residual Analysis}
            \paragraph{가속도계 검증}
                정지 상태에서 보정된 가속도계 출력의 크기가 중력 가속도 $g$와 일치하는지 확인한다:\\

                \begin{equation}
                    \epsilon_a = \left\| \mathbf{a}_{\text{calib}} \right\| - g
                    \label{eq:accel_residual}
                \end{equation}

                잔차 $\epsilon_a$의 평균 및 표준 편차가 충분히 작으면 캘리브레이션이 유효하다. 6개의 정치 위치 각각에서 잔차를 평가하면 위치 의존적 오차를 식별할 수 있다.\\

            \paragraph{자이로스코프 검증}
                정지 상태에서 바이어스 보정 후 각속도 출력의 평균이 $\mathbf{0}$에 가까운지 확인한다:\\

                \begin{equation}
                    \bar{\boldsymbol{\omega}}_{\text{calib}} = \frac{1}{N}\sum_{k=1}^N \boldsymbol{\omega}_{\text{calib},k} \approx \mathbf{0}
                    \label{eq:gyro_residual}
                \end{equation}

                Allan Variance 재분석을 통해 보정 후 잡음 특성이 개선되었는지 확인한다.\\

            \paragraph{지자기계 검증}
                다양한 방향에서 보정된 자기장 크기의 표준 편차를 계산한다:\\

                \begin{equation}
                    \epsilon_m = \text{std}\left( \left\| \mathbf{m}_{\text{calib}} \right\| \right)
                    \label{eq:mag_residual}
                \end{equation}

                보정 전 $\epsilon_m$이 크게 감소하면 캘리브레이션이 유효하다. 이상적으로는 $\epsilon_m / \|\mathbf{m}_{\text{true}}\| < 1\%$를 목표로 한다.\\

        \subsubsection{Residual Analysis with Static Test}
            외부 기준 센서가 없는 환경에서는 자체 일관성(self-consistency) 검증을 수행한다.\\

            \paragraph{정적 자세 반복성 검증}
                동일한 정치 위치에서 여러 번 자세를 추정하고 반복성을 평가한다. 가속도계로부터 추정한 roll/pitch 각도의 표준 편차가 캘리브레이션 후 감소해야 한다:\\

                \begin{equation}
                    \sigma_{\phi,\text{after}} < \sigma_{\phi,\text{before}}, \qquad \sigma_{\theta,\text{after}} < \sigma_{\theta,\text{before}}
                    \label{eq:static_repeatability}
                \end{equation}

            \paragraph{온도 안정성 검증}
                가능하다면 다양한 온도 조건에서 바이어스 추정값의 변화를 측정하여 온도 계수(temperature coefficient)를 평가한다. 온도 보정은 본 연구의 범위를 벗어나나, 온도 변화가 클 경우 주기적 재캘리브레이션이 필요함을 인지해야 한다.\\

            \paragraph{캘리브레이션 파라미터 요약}
                캘리브레이션 결과는 표 \ref{tab:calib_params}와 같이 정리하여 5장의 EKF 초기화 및 Q/R 설정에 활용한다.\\

                \begin{table}[htbp!]
                    \centering
                    \caption{캘리브레이션 파라미터 요약}
                    \label{tab:calib_params}
                    \begin{tabular}{llll}
                        \toprule
                        센서 & 파라미터 & 기호 & 활용 \\
                        \midrule
                        가속도계 & 바이어스 & $\mathbf{b}_a$ & 측정값 보정 \\
                                & 스케일 팩터 & $\mathbf{S}_a$ & 측정값 보정 \\
                                & 잡음 표준편차 & $\sigma_a$ & EKF $\mathbf{R}$ \\
                        \midrule
                        자이로스코프 & 정적 바이어스 & $\hat{\mathbf{b}}_g$ & 측정값 보정 \\
                                    & ARW 계수 & $N$ & EKF $\mathbf{Q}$ ($\sigma_w^2$) \\
                                    & RRW 계수 & $K$ & EKF $\mathbf{Q}$ ($\sigma_{rw}^2$) \\
                        \midrule
                        지자기계 & Hard Iron 바이어스 & $\hat{\mathbf{b}}_m$ & 측정값 보정 \\
                                & Soft Iron 행렬 & $\mathbf{W}^{-1}$ & 측정값 보정 \\
                                & 잡음 표준편차 & $\sigma_m$ & EKF $\mathbf{R}$ \\
                        \bottomrule
                    \end{tabular}
                \end{table}